\documentclass[10pt,twocolumn]{article}

% ----------------------------
% Packages
% ----------------------------
\usepackage[margin=1in]{geometry}
\usepackage{times}
\usepackage{amsmath, amssymb}
\usepackage{graphicx}
\usepackage{subcaption}
\usepackage{booktabs}
\usepackage{hyperref}
\usepackage{microtype}

% ----------------------------
% Title + Authors
% ----------------------------
\title{
Do Latent World Models Exhibit Physical Expectation?\\
\large A Controlled Study of Representation Geometry and Predictive Structure
}

\author{
Daniel Enriquez \\
\texttt{dee4002@nyu.edu}
\and
Sihan Ma \\
\texttt{sm11627@nyu.edu}
\and
Ayush Sarkar \\
\texttt{as17812@nyu.edu}
}

\date{}

% ----------------------------
% Document
% ----------------------------
\begin{document}

\maketitle

% ----------------------------
% Abstract
% ----------------------------
\begin{abstract}
A central goal of both cognitive science and machine learning is to understand whether learned representations support predictive internal models of the world. Recent work on world models, particularly predictive latent architectures, suggests that such systems may learn structured representations that enable reasoning about future states. A common evaluation paradigm, inspired by violation-of-expectation (VoE) experiments in developmental psychology, interprets increased prediction error under physically implausible outcomes as evidence of expectation-like behavior.

In this work, we critically examine this assumption. We evaluate representations learned under predictive, contrastive, and reconstruction-based objectives across multiple environments, focusing not on how they are learned, but on what they encode after training. Using a suite of probing, geometric, and prediction-based experiments, we find that naive VoE signals are largely explained by latent geometry: prediction error increases with distance in representation space rather than with structured violations of physical plausibility. When latent distance is controlled, the apparent violation signal collapses. Furthermore, prediction error is largely invariant to action perturbations, indicating weak action-conditioned dynamics.

These results suggest that current VoE-style evaluations conflate smooth representational structure with predictive expectation. From a cognitive perspective, this highlights a distinction between representations that encode relational similarity and those that support causal, action-conditioned prediction. From a machine learning perspective, our findings call for more rigorous evaluation protocols for assessing world models and predictive representations.
\end{abstract}

% ----------------------------
% 1. Introduction
% ----------------------------
\section{Introduction}

A central hypothesis in cognitive science is that intelligent systems rely on internal models of the world that support prediction and action. In Craik's formulation, thought consists of constructing internal representations that can be manipulated to anticipate the consequences of external processes \citep{craik1944nature}. These representations are not merely descriptive: they must support transformations that mirror the causal structure of the environment. Similarly, Tolman argued that behavior is guided by structured internal representations---``cognitive maps''---rather than simple stimulus--response associations \citep{tolman1948cognitive}. In both accounts, intelligence depends not only on encoding the world, but on representing how the world evolves under action.

Recent work in machine learning has revisited these ideas through the development of \emph{world models}: learned systems that aim to represent and predict the dynamics of environments. In particular, a growing line of research, including Joint Embedding Predictive Architectures (JEPA) and related approaches, emphasizes learning representations that are predictive rather than reconstructive \citep{lecun2023jepa, lecun2024worldmodels}. These models aim to capture the structure of future states without explicitly reconstructing raw sensory input, focusing instead on learning latent representations that support prediction over relevant aspects of the world.

This predictive paradigm contrasts with earlier representation learning approaches such as reconstruction-based methods (e.g., masked autoencoders) and contrastive methods. Reconstruction-based models are optimized to reproduce input data at the pixel level, encouraging representations that preserve detailed sensory information but may not prioritize task-relevant or predictive structure. As argued in recent work, such objectives can lead to representations that encode information about the world without capturing the underlying dynamics that govern its evolution \citep{lecun2023jepa}. In contrast, predictive world models aim to encode abstractions that are stable, compressive, and useful for anticipating future states.

Despite this shift in learning objectives, evaluating whether learned representations actually exhibit predictive structure remains challenging. A common approach, inspired by cognitive science, is to use violation-of-expectation (VoE) paradigms. In developmental psychology, VoE experiments infer the presence of internal expectations by measuring increased surprise when observed outcomes violate physical constraints. In machine learning, analogous evaluations measure prediction error under mismatched or counterfactual futures, interpreting higher error as evidence that the model has learned to expect physically plausible outcomes.

However, it is unclear whether such signals genuinely reflect predictive, action-conditioned expectations, or whether they arise from simpler properties of representation space. In particular, latent representations learned by neural networks often exhibit smooth geometric structure, where distances reflect similarity between observations. As a result, prediction error may increase under counterfactual targets simply because they are farther away in latent space, rather than because they violate a learned model of physical dynamics.

In this paper, we investigate this question directly. We study representations learned under different training paradigms---including predictive world models, contrastive encoders, and reconstruction-based models---and evaluate what they encode \emph{after training}. Rather than focusing on the learning process itself, we ask a more targeted question:

\begin{quote}
\emph{Once a representation has been learned, does it support structured, action-conditioned prediction of future states, or does it primarily reflect geometric similarity between observations?}
\end{quote}

To answer this, we introduce a controlled experimental framework that combines probing, geometric analysis, and violation-of-expectation tests with a set of critical controls. We show that while naive VoE metrics suggest strong sensitivity to physical violations, these effects can be explained by latent geometry alone. When latent distance is controlled, the violation signal collapses. Furthermore, prediction error is largely invariant to action perturbations, indicating that the representations do not strongly encode action-conditioned dynamics.

These findings suggest that current evaluation paradigms may overinterpret geometric structure as evidence of predictive understanding. From a cognitive perspective, this distinction parallels Craik's critique of purely relational theories of explanation: systems may capture statistical relationships without modeling the causal processes that generate them. From a machine learning perspective, our results highlight the need for evaluation methods that distinguish between representations that encode the world and those that encode how the world changes.

\paragraph{Contributions.}
We make the following contributions:
\begin{itemize}
    \item We provide a controlled evaluation framework for analyzing predictive structure in learned representations across multiple environments.
    \item We show that naive violation-of-expectation signals are largely explained by latent geometry rather than structured prediction.
    \item We introduce control experiments that isolate the roles of latent distance and action dependence in prediction error.
    \item We demonstrate that reconstruction and contrastive models can exhibit apparent VoE effects without encoding action-conditioned dynamics.
\end{itemize}

% ----------------------------
% 2. Background
% ----------------------------
\section{Background}

\subsection{Predictive World Models and Representation Learning}

Modern approaches to world modeling emphasize learning latent representations that support prediction over future states. In predictive architectures such as JEPA, models are trained to map observations into a latent space and predict future representations conditioned on context and action, without reconstructing the full sensory input \citep{lecun2023jepa}. This framework aims to capture the invariant, task-relevant structure of the world while discarding unpredictable or irrelevant details.

This contrasts with reconstruction-based methods, such as masked autoencoders, which optimize for pixel-level fidelity. While such models can learn rich representations of visual structure, their objectives do not explicitly encourage capturing the dynamics of the environment. Similarly, contrastive methods learn representations that preserve semantic similarity but do not directly model temporal evolution.

A central question, therefore, is whether predictive training objectives lead to representations that differ qualitatively from reconstruction-based or contrastive approaches, particularly in their ability to support prediction and reasoning about future states.

\subsection{Prediction, Causality, and Internal Models}

Craik's theory of thought frames prediction as the manipulation of internal symbols that correspond to external processes \citep{craik1944nature}. For a representation to support prediction in this sense, it must encode not only the current state of the world but also the transformations induced by actions. This aligns with causal theories of explanation, in which understanding involves modeling how interventions produce changes in outcomes.

In contrast, purely relational or descriptive representations capture statistical associations between observations without modeling the underlying processes. Such representations may support similarity judgments or interpolation but do not necessarily enable structured prediction.

\subsection{Cognitive Maps and Structured Representation}

Tolman's concept of cognitive maps provides a complementary perspective, emphasizing structured internal representations that support flexible navigation and goal-directed behavior \citep{tolman1948cognitive}. Cognitive maps encode relationships between states in a way that allows agents to reason about paths, goals, and counterfactual scenarios. In modern terms, this suggests a representation that organizes the environment in a way that reflects its underlying structure, rather than merely embedding observations in a smooth space.

\subsection{Evaluating Representations via Alignment and Expectation}

A common approach to evaluating learned representations is to measure their alignment with human similarity judgments or neural data \citep{sucholutsky2023alignment, muttenthaler2023improving}. Techniques such as multidimensional scaling embed stimuli into a latent space where distances reflect perceived similarity. While these methods provide insight into the structure of representations, alignment alone does not imply that a system supports the same computations or behaviors as humans.

Violation-of-expectation paradigms extend this idea by probing predictive structure. In both cognitive science and machine learning, increased surprise under unexpected outcomes is interpreted as evidence of internal expectations. However, this interpretation depends on the assumption that prediction error reflects structured, action-conditioned models of the world, rather than generic properties of representation geometry.

% ----------------------------
% 2. Method
% ----------------------------
\section{Method}

\subsection{Environments}

We evaluate representations across three environments that vary in their dynamical structure, observability, and task complexity: \textbf{PushT}, \textbf{TwoRoom}, and \textbf{Reacher}. These environments are chosen to probe different aspects of representation, including object interaction, spatial navigation, and continuous control.

\paragraph{PushT.}
PushT is a planar manipulation environment in which an agent must move a block to a target location through physical interaction. The state includes both agent and object properties, such as position, orientation, and velocity. This environment introduces non-trivial dynamics arising from contact and interaction, requiring representations to capture relationships between multiple entities. From a cognitive perspective, PushT tests whether representations encode structured object-level relationships and support prediction in interactive settings.

\paragraph{TwoRoom.}
TwoRoom is a discrete navigation environment in which an agent moves between spatial regions separated by a barrier. The state is defined by the agent's position, and success depends on reaching a goal location that may lie in a different room. This environment emphasizes spatial structure and connectivity, aligning closely with classical cognitive map paradigms \citep{tolman1948cognitive}. Because the underlying state is low-dimensional and fully observable, TwoRoom provides a setting in which representation quality can be cleanly interpreted.

\paragraph{Reacher.}
Reacher is a continuous control task involving a two-link robotic arm that must reach a target position. Observations include joint configurations and task-relevant features such as fingertip and target positions. Unlike PushT, the dynamics are governed by smooth, continuous control rather than contact interactions. This environment tests whether representations capture kinematic structure and support prediction in continuous dynamical systems.

\paragraph{Summary.}
Together, these environments span a range of representational demands: structured object interaction (PushT), spatial abstraction (TwoRoom), and continuous control (Reacher). This diversity allows us to evaluate whether observed effects generalize across settings with different underlying structure, rather than arising from a specific environment or task.

\subsection{Representations}

We evaluate four classes of representations corresponding to distinct training paradigms: predictive world models, contrastive encoders, reconstruction-based models, and random features. These choices allow us to disentangle the role of training objective in shaping the structure of learned representations.

\paragraph{Predictive World Model (LeWM).}
We use latent representations from a predictive world model trained using a joint-embedding predictive objective in the spirit of JEPA \citep{lecun2023jepa}. The model maps observations to a latent space and is trained to predict future latent states conditioned on context and actions, without reconstructing raw sensory inputs. This objective encourages representations that capture invariant, task-relevant structure and support prediction over future states. In contrast to reconstruction-based methods, predictive world models explicitly aim to encode the dynamics of the environment and are therefore the closest analogue to internal models in the sense of \citet{craik1944nature}.

\paragraph{Contrastive Representation (DINOv2).}
We include representations from a contrastive vision model trained using self-distillation (DINOv2). Such models are optimized to produce embeddings that are invariant to augmentations while preserving semantic similarity. While these representations are highly effective for recognition and alignment tasks, they are not trained to model temporal dynamics or predict future states. As a result, they provide a useful baseline for assessing whether predictive structure emerges without explicit predictive supervision.

\paragraph{Reconstruction-Based Representation (MAE).}
We use representations from a masked autoencoder (MAE), trained to reconstruct missing portions of input images. Reconstruction-based objectives encourage the model to preserve detailed sensory information, including high-frequency visual features. However, as emphasized in prior work \citep{lecun2023jepa}, such objectives do not explicitly prioritize learning the dynamics of the environment. In particular, accurate reconstruction does not require encoding how states evolve under action, and thus MAE representations should not be considered world models in the predictive sense.

\paragraph{Random Features.}
As a control, we include randomly initialized features with matched dimensionality. These features provide a baseline for evaluating whether observed effects arise from learned structure or from generic properties of high-dimensional embeddings.

\paragraph{Summary.}
These representations span a spectrum of training objectives, from explicitly predictive (LeWM) to purely reconstructive (MAE) and invariant (DINOv2). By comparing their behavior across tasks, we aim to isolate which aspects of representation---predictive structure, semantic similarity, or raw information preservation---contribute to observed effects in probing, geometry, and violation-of-expectation evaluations.

\subsection{Experimental Framework}

We evaluate representations using a sequence of experiments that isolate what is encoded, how it is organized, and whether it supports predictive expectations. Let $x_t$ denote observations, $a_t$ actions, and $z_t = \phi(x_t)$ the latent representation produced by a given encoder $\phi$.

\paragraph{Experiment 1: Representation Probing.}
We first assess what information is linearly decodable from each representation. For each environment, we train linear probes $g$ to predict ground-truth state variables $s_t$ from latent embeddings:
\begin{equation}
\hat{s}_t = g(z_t).
\end{equation}
This experiment establishes whether representations encode task-relevant physical variables (e.g., positions, orientations, kinematic features). However, high decodability does not imply that representations support structured prediction.

\paragraph{Experiment 2: Latent Geometry.}
Next, we evaluate how representations organize physical state in latent space. Let $d_{\text{lat}}(z_i, z_j)$ denote Euclidean distance in latent space, and $d_{\text{phys}}(s_i, s_j)$ a task-specific physical distance. We measure alignment via rank correlation:
\begin{equation}
\rho = \text{Spearman}(d_{\text{lat}}(z_i, z_j),\; d_{\text{phys}}(s_i, s_j)).
\end{equation}
This tests whether latent spaces behave like structured embeddings of the environment, analogous to cognitive maps.

\paragraph{Experiment 3: Violation-of-Expectation (VoE).}
We evaluate predictive structure by training a transition predictor $f$ in latent space. Given a history of latent states and actions, we predict a future latent:
\begin{equation}
\hat{z}_{t+h} = f\big(z_{t-k:t},\; a_{t:t+h}\big).
\end{equation}

We define prediction error with respect to a target latent $z_{\text{target}}$:
\begin{equation}
S = \|\hat{z}_{t+h} - z_{\text{target}}\|_2.
\end{equation}

We evaluate four conditions:
\begin{align*}
\text{normal:} &\quad z_{\text{target}} = z_{t+h} \\
\text{near:} &\quad z_{\text{target}} \sim \mathcal{N}_{\text{near}}(s_{t+h}) \\
\text{far:} &\quad z_{\text{target}} \sim \mathcal{N}_{\text{far}}(s_{t+h}) \\
\text{random:} &\quad z_{\text{target}} \sim \text{Uniform}(\mathcal{D})
\end{align*}
where $\mathcal{N}_{\text{near}}$ and $\mathcal{N}_{\text{far}}$ denote sets of states selected by percentile ranges of physical distance from the true future.

A standard interpretation is that
\begin{equation}
S_{\text{normal}} < S_{\text{near}} < S_{\text{far}}
\end{equation}
indicates expectation-like predictive structure.

\paragraph{Control 1: Latent-Distance Matching.}
To test whether this effect reflects physical plausibility or latent geometry, we introduce a control that matches targets by latent distance. For each $(z_{t+h}, z_{\text{target}})$ pair, we select comparison targets $\tilde{z}$ such that:
\begin{equation}
\|z_{t+h} - \tilde{z}\|_2 \approx \|z_{t+h} - z_{\text{target}}\|_2,
\end{equation}
while varying physical distance. If prediction error is driven by physical structure, the violation effect should persist; if driven by geometry, it should collapse.

\paragraph{Control 2: Action Perturbation.}
We assess action dependence by replacing the true action sequence with a shuffled sequence $\tilde{a}_{t:t+h}$:
\begin{equation}
\hat{z}_{t+h}^{\,\text{shuffle}} = f\big(z_{t-k:t},\; \tilde{a}_{t:t+h}\big).
\end{equation}

We then compute the action sensitivity ratio:
\begin{equation}
R_{\text{action}} = \frac{\mathbb{E}[\|\hat{z}_{t+h}^{\,\text{shuffle}} - z_{t+h}\|_2]}{\mathbb{E}[\|\hat{z}_{t+h} - z_{t+h}\|_2]}.
\end{equation}

Values near $1$ indicate weak dependence on action sequences, suggesting that prediction is driven primarily by state similarity rather than action-conditioned dynamics.

\paragraph{Summary.}
Together, these experiments separate three questions: (1) what representations encode, (2) how they structure physical state, and (3) whether they support action-conditioned prediction. The control conditions allow us to distinguish genuine predictive structure from effects arising from smooth latent geometry.

% ----------------------------
% 3. Results
% ----------------------------
\section{Results}

\subsection{Representation Geometry}

\begin{figure}[t]
    \centering
    \includegraphics[width=\linewidth]{paper_fig1_geometry.pdf}
    \caption{Latent–physical geometry alignment across environments.}
    \label{fig:geometry}
\end{figure}

% text later

\subsection{Naive Violation-of-Expectation}

\begin{figure}[t]
    \centering
    \includegraphics[width=\linewidth]{paper_fig2_naive_voe.pdf}
    \caption{Naive VoE effect: prediction error increases under physically distant targets.}
    \label{fig:naive_voe}
\end{figure}

% text later

\subsection{Latent-Distance Control}

\begin{figure}[t]
    \centering
    \includegraphics[width=\linewidth]{paper_fig3_latent_control.pdf}
    \caption{VoE effect under latent-distance matching. The effect collapses across models.}
    \label{fig:latent_control}
\end{figure}

% text later

\subsection{Action Dependence}

\begin{figure}[t]
    \centering
    \includegraphics[width=\linewidth]{paper_fig4_action_dependence.pdf}
    \caption{Effect of action shuffling on prediction error. Ratios near 1 indicate weak action dependence.}
    \label{fig:action_control}
\end{figure}

% text later

\subsection{Summary of Controls}

\begin{figure}[t]
    \centering
    \includegraphics[width=\linewidth]{paper_fig5_summary.pdf}
    \caption{Naive vs controlled VoE effect. Points collapse toward 1 under controls.}
    \label{fig:summary}
\end{figure}

% text later

% ----------------------------
% 4. Discussion
% ----------------------------
\section{Discussion}

% (we will write this next, section by section)

\subsection{Apparent Violation-of-Expectation as a Geometric Artifact}

% text later

\subsection{Limits of Action-Conditioned Prediction}

% text later

\subsection{Implications for Evaluating World Models}

% text later

% ----------------------------
% 5. Conclusion
% ----------------------------
\section{Conclusion}

% (short)

% ----------------------------
% References
% ----------------------------
\bibliographystyle{plain}
\bibliography{references}

\end{document}